\documentclass[11pt, headheight=30pt]{scrartcl}
\usepackage[white, nologo, hw]{darkWhite}

\DeclareMathOperator{\Aut}{Aut}
\DeclareMathOperator{\Der}{Der}
\DeclareMathOperator{\Rad}{Rad}

% sp
\DeclareMathOperator{\spp}{\mathfrak{sp}}

\author{Rowechen Zhong (Collaborator: Isaac Zhu)}
\title{18.745 Lie Algebra}
\subtitle{Homework 7}
\begin{document}
% 18.745: Homework 7
% Due 11:59pm Wednesday Dec 13
% In the following, let L be a semisimple Lie algebra. Let Φ be a root system with respect
% to a CSA H. Fix a base ∆ = {α1, · · · , αℓ} and n+ =
% P
% α∈Φ+ Lα.

% 1. Let {λ1, · · · , λℓ} be the fundamental dominant weights. Show how to construct an arbitrary V (λ), λ ∈ Λ
% + as a direct summand in a suitable tensor product of V (λ1), · · · , V (λℓ).
\section{Representation Decomposition}

In the following, let $L$ be a semisimple Lie algebra. Let $\Phi$ be a root system with respect to a CSA $H$. Fix a base $\Delta = \{\alpha_1, \dots, \alpha_\ell\}$ and $\mathfrak{n}_+ = \sum_{\alpha \in \Phi_+} L_\alpha$.

\begin{prob*}
Let $\{\lambda_1, \dots, \lambda_\ell\}$ be the fundamental dominant weights. Show how to construct an arbitrary $V(\lambda), \lambda \in \Lambda_+$ as a direct summand in a suitable tensor product of $V(\lambda_1), \dots, V(\lambda_\ell)$.
\end{prob*}

This is pretty clear. Suppose $\lambda = \sum a_i \lambda_i$. Then let $V = V(\lambda_1)^{\otimes a_1} \otimes \cdots \otimes V(\lambda_\ell)^{\otimes a_\ell}$. The highest weight vector of $V$ is $v = v_1 \otimes \cdots \otimes v_\ell$ where $v_i$ is the highest weight vector of $V(\lambda_i)$; the weight is exactly
\[
    \lambda_1 + \cdots + \lambda_\ell = \sum a_i \lambda_i = \lambda.    
\]
(this is clear, because if you act $v$ with an element $h$ of the CSA, then $h$ acts on each $v_i$ by $\lambda_i(h)$, so the total weight is $\sum \lambda_i(h) = \lambda(h)$). Thus $V(\lambda)$ is a direct summand of $V$.

% 2. Let 1 ≤ k ≤ n. Consider the k-th exterior power representation W = ΛkC
% n of sln(C).
% Let {⃗e1, · · · , ⃗en} be the standard basis of C
% n
% . Let H be the standard CSA (which consists
% of diagonal matrices) and take the standard ∆ ⊂ Φ
% + ⊂ Φ.
% (1) Compute all weights of W and describe the corresponding weight spaces.
% (2) Show that W contains a unique vector ⃗w such that n+ ⃗w = 0, namely ⃗v = ⃗e1 ∧⃗e2 ∧
% · · · ∧ ⃗ek, and deduce from this that W is irreducible.
% (3) Find the highest weight of W.

\section{Exterior Power Representation}
\begin{prob*}
Let $1 \leq k \leq n$. Consider the $k$-th exterior power representation $W = \Lambda^k \CC^n$ of $\sll(n, \CC)$. Let $\{\vec{e}_1, \dots, \vec{e}_n\}$ be the standard basis of $\CC^n$. Let $H$ be the standard CSA (which consists of diagonal matrices) and take the standard $\Delta \subset \Phi_+ \subset \Phi$.
\begin{enumerate}
    \item Compute all weights of $W$ and describe the corresponding weight spaces.
    \item Show that $W$ contains a unique vector $\vec{w}$ such that $\mathfrak{n}_+ \vec{w} = 0$, namely $\vec{v} = \vec{e}_1 \wedge \vec{e}_2 \wedge \cdots \wedge \vec{e}_k$, and deduce from this that $W$ is irreducible.
    \item Find the highest weight of $W$.
\end{enumerate}
\end{prob*}

Let $L = \sll(n, \CC)$, the traceless matrices. The root space lives in $E$, $\RR^n$ modulo $\sum e_i = 0$. The root basis is $\{e_i - e_{i+1}\}$ corresponding to diagonal elements $(M^i)_{jk} = \delta_{jk}\left(\delta_{j,i} - \delta_{j,i+1}\right)$ for $1 \leq i < n$.

The weights of basis vector $\vec{e}_i\in \CC^n$ is precisely $e_i\in E$ (using different letters to avoid confusion). A basis for $W$ is
\[
    \vec{e}_i \equiv \vec{e}_{i_1} \wedge \cdots \wedge \vec{e}_{i_k}    
\]
with $1 \leq i_1 < \cdots < i_k \leq n$. The weight of this basis vector is exactly $e_{i_1} + \cdots + e_{i_k}$. Thus the weights of $W$ are exactly the sums of $k$-element subsets of $\{\vec{e}_1, \dots, \vec{e}_n\}$. Each weight space has dimension one and contains exactly the vector above.


Next, observe $\mathfrak{n}_+ v = 0$ iff $L_\alpha v = 0$ for all $\alpha \in \Phi_+$. We can explicitly write the positive root spaces in our root system; under the basis above, they are exactly the matrices
\[
    (A^{ij})_{kl} = \delta_{k,i}\delta_{l,j}
\]
with $i < j$. The basis is precisely when $j = i + 1$. We will now consider a general element expressed as
\[
    \sum_{i} c_i \vec{e}_{i}  
\]
where the index $i$ ranges over all $k$-element increasing subsequences of $\{1, \dots, n\}$. These subsets can be partially ordered by $i \succ j$ iff there exists $k$ with $i_k > j_k$. Suppose $\mathfrak{n}_+ v = 0$, but $v\neq 0$. Pick any maximimal $i$ with nonzero $c_i$. Suppose $i\neq (1, \dots, k)$.
Then,
\[
    A^{ki_k} A^{(k-1)(i_k - 1)}\dots A^{i_1 1}\left(\sum_{i} c_i \vec{e}_i\right) = 0.
\]
because the sequence of $A$'s is nonempty because $i\neq (1,\dots, k)$.
However, observe that the LHS annihilates any $\vec{e}_j$ for $j\neq i$ by inspection. Thus the LHS is exactly $c_i \vec{e}_{(1,\dots k)}\neq 0$. Thus, $A$ must be empty, and we have the desired.

The highest weight is just $f_1 + \cdots + f_k$.

% 3. Let λ ∈ Λ
% +. Prove that the dimension of the irreducible module V (λ) is given by
% dim V (λ) = Y
% α∈Φ+
% ⟨λ + δ, α⟩
% ⟨δ, α⟩
\section{Dimension of Irreducible Modules}
\begin{prob*}
Let $\lambda \in \Lambda_+$. Prove that the dimension of the irreducible module $V(\lambda)$ is given by
\[
    \dim V(\lambda) = \prod_{\alpha \in \Phi_+} \frac{\langle \lambda + \delta, \alpha \rangle}{\langle \delta, \alpha \rangle}
\]
\end{prob*}

This is the corollary on page $139$ of Humphrey's; the proof is provided there also. The basic idea is to cite Harish-Chandra, leading to Kostant's multiplicities theorem. Then we convolve to get Weyl's theorem, which quickly yields the desired result. Given the amount of information it is necessary to cite, I will not reproduce the proof here.

% 4. Let L = sp6
% (1) Consider the adjoint representation ad : sp6 → gl(L). Prove that it is an irreducible
% representation.
% (2) Find the maximal weight and describe a maximal vector in L.
\section{Adjoint Representation of $\spp(6, \CC)$}
\begin{prob*}
Let $L = \spp(6, \CC)$.
\begin{enumerate}
    \item Consider the adjoint representation $\ad : \spp(6, \CC) \to \gl(L)$. Prove that it is an irreducible representation.
    \item Find the maximal weight and describe a maximal vector in $L$.
\end{enumerate}
\end{prob*}

The weight space of the adjoint representation is exactly the root space plus the origin (corresponding to the CSA). This root space is exactly $C_3$, but this Dynkin diagram is connected, so the adjoint representation is irreducible. (This proof shows all simple Lie algebras have irreducible adjoint representations.)

For the second part, we compute some things specific to $\spp(2n, \CC)$, $n = 3$. This consists of the block matrices
\[
    \begin{pmatrix}
        A & B \\
        C & -A^T
    \end{pmatrix}
\]
where $B$ and $C$ are symmetric; it is $n(2n + 1) = 21$-dimensional.
% The CSA consists of the diagonal matrices
% \[
%     \begin{pmatrix}
%         D & 0 \\
%         0 & -D
%     \end{pmatrix}
% \]
% where $D$ is diagonal; it is $n = 3$ dimensional.

We now write out the CSA and root space decomposition, which is essentially a trivial but tedious computation. In the below, I will write out only the general cases. The CSA consists of
\[
    H_1 = \left(
        \begin{array}{ccc|ccc}
            1 & 0 & 0 & 0 & 0 & 0 \\
            0 & 0 & 0 & 0 & 0 & 0 \\
            0 & 0 & 0 & 0 & 0 & 0 \\
            \hline
            0 & 0 & 0 & -1 & 0 & 0 \\
            0 & 0 & 0 & 0 & 0 & 0 \\
            0 & 0 & 0 & 0 & 0 & 0    
        \end{array}
    \right), \dots
\]
(3 total). The root spaces of the form $(1,-1,0)$ look like
\[
    X_{1,2} = \left(
        \begin{array}{ccc|ccc}
            0 & 1 & 0 & 0 & 0 & 0 \\
            0 & 0 & 0 & 0 & 0 & 0 \\
            0 & 0 & 0 & 0 & 0 & 0 \\
            \hline
            0 & 0 & 0 & 0 & 0 & 0 \\
            0 & 0 & 0 & -1 & 0 & 0 \\
            0 & 0 & 0 & 0 & 0 & 0    
        \end{array}
    \right), \dots
\]
(6 total). The root spaces of the form $(1,1,0)$ look like
\[
    Y_{1,2} = \left(
        \begin{array}{ccc|ccc}
            0 & 0 & 0 & 0 & 1 & 0 \\
            0 & 0 & 0 & 1 & 0 & 0 \\
            0 & 0 & 0 & 0 & 0 & 0 \\
            \hline
            0 & 0 & 0 & 0 & 0 & 0 \\
            0 & 0 & 0 & 0 & 0 & 0 \\
            0 & 0 & 0 & 0 & 0 & 0    
        \end{array}
    \right)
\]
(3 total, and their transposes $Y^\top_{i,j}$ yield the $(-1,-1,0)$ root spaces). The root spaces of the form $(2,0,0)$ look like
\[
    Z_1 = \left(
        \begin{array}{ccc|ccc}
            0 & 0 & 0 & 1 & 0 & 0 \\
            0 & 0 & 0 & 0 & 0 & 0 \\
            0 & 0 & 0 & 0 & 0 & 0 \\
            \hline
            0 & 0 & 0 & 0 & 0 & 0 \\
            0 & 0 & 0 & 0 & 0 & 0 \\
            0 & 0 & 0 & 0 & 0 & 0
        \end{array}
    \right), \dots
\]
(3 total, and their transposes $Z^\top_i$ yield the $(-2,0,0)$ root spaces).

Thus the root system is exactly $\{\pm 2e_i, \pm e_i \pm e_j\}$ for $1\leq i < j\leq n = 3$. The Cartan matrix of $C_3$ is
\[
    \begin{pmatrix}
        2 & -1 & 0 \\
        -1 & 2 & -1 \\
        0 & -2 & 2
    \end{pmatrix},   
\]
and the root system looks like
\begin{center}
    \includegraphics[width=0.5\textwidth]{c3.png}
\end{center}
We pick a basis $\{e_1 - e_2, e_2 - e_3, 2e_3\}$ (shown in red) of $L$. Then, a maximal weight (which is the same as a maximal root) is $\alpha = 2e_1$ (shown in blue); it is the unique maximum under the partial order $\alpha \leq \beta$ if $\beta - \alpha$ is a nonnegative linear combination of positive roots. It corresponds to a block of the form $Z_i$.

% 5. Suppose L is simple. Show that the universal Casimir element cL is in the center Z of
% the universal enveloping algebra U(L).
% 1
\section{Universal Casimir Element}
\begin{prob*}
Suppose $L$ is simple. Show that the universal Casimir element $c_L$ is in the center $Z$ of the universal enveloping algebra $U(L)$.
\end{prob*}

Recall that the universal Casimir element is
\[
    c_L = \sum_{i=1}^n x_i\otimes y_i \in U(L),
\]
where $\{x_i\}$ is a basis of $L$ and $\{y_i\}$ is the dual basis under the Killing form. Compute:
\begin{align*}
    [c_L, z]
    &= \sum_{i=1}^n [x_i \otimes y_i, z]\\
    &= \sum_{i=1}^n x_i\otimes [y_i, z] + [x_i, z]\otimes y_i
\end{align*}
Contract the second tensor element against any $x_j$ using $B$.
\begin{align*}
    &= \sum_{i=1}^n x_i\otimes B([y_i, z], x_j) + [x_i, z]\otimes B(y_i, x_j)\\
    &= \sum_{i=1}^n x_i\otimes B(y_i, [z, x_j]) + [x_j, z]
\end{align*}
Letting $w = [x_j, z]$, we observe that this is true for all $w = x_k$ thus we conclude.

\end{document}